\section{Samarbeidsmodell, Totalforsvar og ROS-analyse}
\label{sec:totalforsvar}

For at systemet skal aktiveres på under to timer, etableres et forankret offentlig-privat samvirke (\textsc{ops}) innenfor rammen av det norske Totalforsvaret.

\subsection{Totalforsvarsmodellen og ansvarsfordeling}
Konseptet forankres i et formalisert grensesnitt mellom tre aktører:
\begin{itemize}[leftmargin=1.5em, itemsep=0.2em]
    \item \textbf{Meieriaktørene (TINE og Q-Meieriene):} Stiller tankbiler, sjåfører og \textsc{cip}-stasjoner til disposisjon via forhåndsavtaler, og omdirigerer biler via sine telematikksentraler.
    \item \textbf{Tilsynsmyndigheter (Mattilsynet og \textsc{fhi}):} Etablerer forhåndsgodkjente kriserutiner. Digital \textsc{cip}-logg og negative hurtigtester (\textsc{atp} og kasein/myse) fungerer som umiddelbar friskmelding.
    \item \textbf{Kommunal kriseledelse, Sivilforsvaret og Røde Kors:} Klargjør reservefyllepunkter og drifter de sivile utdelingsstasjonene med kanistedistribusjon, desinfeksjonssluse og køavvikling.
\end{itemize}

\subsection{Avbøting av konflikt med melkeinnsamling og dyrevelferd}
Melkekyr lakterer kontinuerlig, og gårdstanker fylles opp på 48 timer. Stans i melkehenting skaper akutte dyrevelferdsproblemer (mastitt), truer nødslakt og medfører melkedumping. Prosjektet løser dette med en to-faset flåtestyringsmodell:
\begin{enumerate}[leftmargin=1.5em, itemsep=0.15em]
    \item \textbf{Fase 1: Akuttfasen (første 12--24 timer):} Behovet på 15 biler dekkes utelukkende ved å mobilisere turnusreserver, biler på verkstedbuffer og overskuddskapasitet fra tilgrensende regioner. Ingen ordinære melkeruter berøres.
    \item \textbf{Fase 2: Langvarig krise ($> 24$ timer):} TINE og Q-Meieriene aktiverer unntaksruter med samkjøring og systematisk overgang til henting annenhver dag (48-timerssyklus) der gårdstankene har tilstrekkelig reservevolum.
\end{enumerate}
Gjennom sivilbeskyttelsesloven har livreddende vannforsyning juridisk forrang, men logistikkmodellen sikrer at dyrevelferden og melkeverdiene ivaretas fullt ut.

\subsection{Risiko- og sårbarhetsanalyse (ROS)}
Systemets operative robusthet og barrierer er evaluert i tabell~\ref{tab:ros}.

\begin{table}[htbp]
\centering
\footnotesize
\caption{Risiko- og sårbarhetsanalyse (\textsc{ros}) for Prosjekt Meierikjeden.}
\label{tab:ros}
\begin{tabularx}{\textwidth}{p{2.6cm}p{3.2cm}ccX}
\toprule
\textbf{Risikohendelse} & \textbf{Konsekvens} & \textbf{S} & \textbf{K} & \textbf{Forebyggende barrierer} \\
\midrule
Konflikt med melkeinnsamling & Gårdstanker overfylles; melkedumping og mastitt. & L & H & To-faset flåtestyring; verkstedreserver; 48t unntaksruter. \\
\addlinespace
Frost i tapperigg og ventiler & Manifold fryser ved kuldegrader; stans i drift. & M & M & Integrert \SI{24}{\volt} (\SI{150}{\watt}) varmekabel; hurtigtømmende bunnventiler. \\
\addlinespace
Krysskontaminering ved sivil tappefront & Patogener fra skitne kanner infiserer tapperigg. & M & H & Tilbakeslagsventiler; berøringsfrie dyser; utdeling av foldekanistere; sprit-sluse. \\
\addlinespace
Svikt i veiinfrastruktur / aksellast & Biler (32--50t) hindres av flomskadde broer/ras. & M & H & Sanntidskobling mot \textsc{vts}; ruting kun på godkjente \textsc{bk}10-korridorer. \\
\addlinespace
Melkeprotein- og allergenrester & Allergisk reaksjon hos sensitive pasienter. & L & K & Alkalisk hydrolyse under \textsc{cip}; to-trinns hurtigtest (\textsc{atp} + lateral-flow, $< \SI{1}{ppm}$). \\
\bottomrule
\multicolumn{5}{l}{\textit{S: Sannsynlighet (L/M/H). K: Konsekvens (M/H/K: Middels, Høy, Kritisk).}}
\end{tabularx}
\end{table}
